\section{Numerical method}
\label{sec:method}

\status{Ready to write: every part described here is implemented and verified.
The subsection on adaptive refinement (\ref{sec:amr}) is the exception --- the
criterion exists and is measured in two dimensions, but has not been ported to
three.}

\subsection{Shared machinery, and the single point of difference}
\label{sec:shared}

This subsection carries the methodological claim of the paper, so it states the
arrangement precisely.

Both methods use the same projection method, the same variable-coefficient
momentum solve, and the same property interpolation. The front-tracking
coupling is a function pointer, null by default, invoked inside the multiphase
step \emph{after} the source terms and \emph{before} the viscosity and density
interpolation. At that point it overwrites the volume fraction from the
geometry of the front, so the properties are built from the new fraction. VOF
reaches the same field at the same point by advecting (\ref{eq:transport}) with
a geometric PLIC reconstruction.

Everything downstream of that point --- $\rho(\mathbf{x})$, $\mu(\mathbf{x})$,
the momentum solve, the pressure projection, the output probes --- is
bit-identical code. The consequence is that a measured difference between the
two is a difference in how the interface is represented, and not in the solver,
the mesh, the time step, the initial geometry, or the instrument.

\subsection{VOF: geometric PLIC}

\missing{PLIC reconstruction and the split advection --- short, since it is the
pre-existing method; cite the reference implementation}

\paragraph{Initial volume fraction.}
The initial fraction is computed by \textbf{sub-cell subdivision}: each
interface cell is split into $16^3$ sub-cells, each resolved by the closed-form
plane--box intersection volume, and the sum is condensed onto the coarse cell.
This matters more than it appears. A one-cell smoothed step --- the obvious
shortcut --- gives a bubble volume error of \num{1.23e-2} at the running
resolution against \num{6.11e-5} for the subdivision, and its measured order of
convergence is erratic (1.59, 2.90, 0.99) because the ramp width is tied to $h$.
Since VOF conserves faithfully whatever mass it is given, an initial-condition
error of that size never recovers.

\subsection{Front-tracking: explicit triangulated surface}

The interface is a closed triangulated surface with explicit connectivity,
built as an icosphere by recursive subdivision of an icosahedron.

\paragraph{Surface tension by line integral.}
The force on each triangle is obtained from the line integral along its edges,
$\mathbf{f} = \sigma(\mathbf{b}-\mathbf{a})\times\mathbf{n}$, distributed to the
two endpoints. The force is $-\sigma\nabla A$ --- tension minimises area ---
and the sign was derived rather than guessed, then checked against the numerical
area gradient.

\paragraph{The Laplace oracle in three dimensions.}
For a sphere the pressure jump is $\Delta p = 2\sigma/R$, \emph{not} $\sigma/R$.
The two-dimensional value is $\sigma/R$, and a 2D implementation that closes to
0.02\% against it would, if the number were carried over, validate the wrong
formula. The distinction is flagged in the header, the harness and the example.

\paragraph{Surgery.}
Split, collapse and flip, with a curved midpoint
\begin{equation}
\mathbf{m} = \frac{\mathbf{a}+\mathbf{b}}{2}
  - \frac{1}{8}\Big[\big((\mathbf{b}-\mathbf{a})\cdot\mathbf{n}_a\big)\mathbf{n}_a
  + \big((\mathbf{a}-\mathbf{b})\cdot\mathbf{n}_b\big)\mathbf{n}_b\Big],
\end{equation}
which is the PN-triangle (Hermite) construction. Surgery keeps the triangle
edge matched to the cell size as the bubble deforms.

\paragraph{Volume fraction from the surface.}
The fraction is obtained from the signed distance and normal to the surface,
then from the closed-form plane--box volume. The sign comes from parity
ray-casting, never from the nearest triangle's normal --- the latter is wrong
for a cell straddling a concavity.

\subsection{Interface extraction for comparison}
\label{sec:mc}

VOF has no interface as an object; the comparison needs one. The isosurface
$f = 1/2$ is extracted by marching cubes on the \emph{dual} grid --- the volume
fraction lives at cell centres, so the lattice nodes are cell centres --- with
vertices \textbf{welded} by lattice edge, which makes the extracted surface
continuous. The alternative, writing the per-cell PLIC planes, gives a surface
that is discontinuous from one cell to the next; that discontinuity is the right
object for auditing the reconstruction and the wrong one for measuring shape.

Both methods therefore emit a triangulated surface in the same file format, and
every shape quantity in this paper is computed from those surfaces by the
\emph{same} routine. The instrument is common; only the surface differs.

\paragraph{On adaptive meshes the extraction runs on the finest level only.}
This is sufficient, and not by economy: the band criterion refines to the finest
level with a floor of five fine cells on each side of the interface, and outside
that band the fraction is exactly 0 or 1. The eight-node stencil therefore never
reaches the seam between levels, which is what normally breaks isosurface
extraction on adaptive meshes.

\subsection{Sub-grid undulation, and the filter}
\label{sec:tsur}

\status{Implemented, verified, and \emph{measured on the production runs}. The
with/without comparison at the reference resolution is pending.}

A Lagrangian surface accumulates error as vertex-scale zig-zag, because the
vertices move independently under the interpolated velocity and nothing
constrains the sub-grid shape. Quantifying it needs a shape-independent
indicator, and the natural one is already computed by the filter itself: for an
edge $(v_1,v_2)$ with edge normal $\mathbf{n}_e$ and midpoint $\mathbf{m}$,
\begin{equation}
h_1 = (\mathbf{v}_1 - \mathbf{m})\cdot\mathbf{n}_e ,
\end{equation}
and because $\mathbf{m}$ is the midpoint, $h_2 = -h_1$ exactly --- so $h_1$ is
the edge's \emph{tilt} relative to the local surface. Normalised by the
half-edge it is dimensionless, and on a smooth surface it is residual whatever
the shape. Measuring radial deviation instead, as the synthetic tests do, would
measure deformation on a deformed bubble.

\begin{table}[htbp]
\centering
\caption{Undulation at the final frame. ``Above'' is the fraction of edges with
indicator above $0.05$; the initial sphere is the oracle.}
\label{tab:undulation}
\small
\begin{adjustbox}{max width=\linewidth}
\pgfplotstabletypeset[
  col sep=space, string type,
  columns/caso/.style={column name={Run}, column type={l}},
  columns/acima/.style={column name={above}, column type={r}},
  columns/inferior/.style={column name={lower half}, column type={r}},
  columns/mediana/.style={column name={median}, column type={r}},
  columns/maximo/.style={column name={max}, column type={r}},
  every head row/.style={before row=\toprule, after row=\midrule},
  every last row/.style={after row=\bottomrule},
]{\dat{ondulacao.dat}}
\end{adjustbox}
\end{table}

Two things in that table. An indicator of $0.997$ means an edge nearly
\emph{perpendicular} to the local surface --- a full zig-zag. And the undulation
is not diffuse: of the 1289 edges above threshold in the coarse Newtonian run,
\textbf{all 1289} lie in the lower half of the bubble. In the worst case, 1765
of 1770.

\subsubsection{The two methods differ in kind, not degree}

\begin{table}[htbp]
\centering
\caption{The indicator over time, same case, both methods. The VOF
\emph{level} is not comparable --- see below.}
\label{tab:undtime}
\small
\pgfplotstabletypeset[
  col sep=space, string type,
  columns/t/.style={column name={$t$}, column type={r}},
  columns/ft/.style={column name={FT above (\%)}, column type={r}},
  columns/ftmax/.style={column name={FT max}, column type={r}},
  columns/vof/.style={column name={VOF above (\%)}, column type={r}},
  columns/vofmax/.style={column name={VOF max}, column type={r}},
  every head row/.style={before row=\toprule, after row=\midrule},
  every last row/.style={after row=\bottomrule},
]{\dat{ondulacao-tempo.dat}}
\end{table}

The VOF \emph{absolute} level means nothing: it is already 27\% at $t=0$, where
the surface is the same sphere front-tracking represents with $0.00\%$. That is
faceting of the marching-cubes extraction, not of the interface. Only the trend
is comparable, and the trends are opposite --- front-tracking accumulates from
zero to 21\% with the maximum reaching $0.997$; VOF stays at $32\pm2\%$ with the
maximum never above $0.37$.

The reason is structural. PLIC reconstructs \emph{one plane per cell}, and a
plane cannot undulate within a cell; the reconstruction is re-derived from the
field every step, so nothing accumulates at the surface level. Put the other
way: \textbf{the PLIC reconstruction is itself a filter}, and a far more
aggressive one than the method below --- it re-imposes a planar interface every
step. The price is the dual of the benefit: it cannot represent anything below
the cell, which is exactly what a cusp is.

So the two methods sit at opposite ends of one trade-off: VOF cannot undulate
and cannot resolve sub-grid curvature; front-tracking can resolve it and
contaminates it with undulation at the same scale. That makes the filter the
difference between the two methods being compared on the cusp, rather than one
of them being disqualified from it.

\subsubsection{The filter, and why it needs a criterion}

BALANCE repositions each vertex along its ring normal to the height that
\emph{reproduces} the star's volume; TSUR puts both ends of an edge at a common
height, weighted by the two star volumes. Both are from the reference
Freeflow implementation~\cite{desousa2004}, where the procedure is called a
``mass-conserving filter''.

That property is measured, not taken on faith: on a sphere with an $8\%$ radial
undulation, three passes drop the deviation by $84\%$ with a volume change of
$0.000\%$, where a Laplacian average smooths slightly \emph{more} and loses
$9.74\%$ of the volume.

\paragraph{But conserving volume does not protect against artificial surface
tension.} Reducing area at fixed volume is precisely what surface tension does.
On a \emph{sphere} the effect is invisible, because the smooth state is already
the minimum-area shape and the filter pushes where the physics also pushes. An
ellipsoid exposes it: at fixed volume, 20 passes take the axis ratio from
$2.0239$ to $1.9072$ and the sphericity from $0.925670$ to $0.931380$,
monotonically, with $\mathrm{d}V/V = 2\times10^{-15}$.

That is $0.00029$ of sphericity per pass. Applied every 20 steps over 1501
steps --- 75 applications --- it would induce $0.021$, which is \textbf{nine
times} the sphericity difference between the two methods that this paper
measures ($0.0023$). Unconditional application would destroy the measurement.

\paragraph{The criterion.} Apply only where there \emph{is} undulation, gated on
the indicator above. There is no trade-off to tune; the plateau is wide.

\begin{table}[htbp]
\centering
\caption{Threshold sweep. The second column must stay far below $0.0023$; the
last must stay high.}
\label{tab:threshold}
\small
\adjustbox{max width=\linewidth}{%
\pgfplotstabletypeset[
  col sep=space, string type,
  columns/limiar/.style={column name={threshold}, column type={r}},
  columns/dpsi/.style={column name={$\Delta\Psi$, 75 passes, smooth ellipsoid}, column type={r}},
  columns/movidos/.style={column name={vertices moved/pass}, column type={r}},
  columns/queda/.style={column name={undulation removed}, column type={r}},
  every head row/.style={before row=\toprule, after row=\midrule},
  every last row/.style={after row=\bottomrule},
]{\dat{filtro-limiar.dat}}
}
\end{table}

At any threshold above $0.02$ the smooth ellipsoid is left \emph{completely
untouched} --- zero vertices moved per pass --- and the undulation still drops
by $82\%$. We use $0.05$, with margin from the edge.

\missing{the production with/without pair at the reference resolution. The
without-filter half is measured (Table~\ref{tab:undtime}); the with-filter run
is pending.}

\subsection{Adaptive refinement}
\label{sec:amr}

\status{Implemented and verified in three dimensions, including the dynamic
remeshing loop. No production run on an adapted mesh yet.}

The criterion is the maximum of two: a geometric band around the interface,
derived from a minimum-fine-cell rule, and a vorticity threshold that captures
the wake. Leaving the fine mesh is deliberately harder than entering it --- a
cell already refined only coarsens once $|\omega|$ falls below \emph{half} the
threshold. Without that hysteresis the two-dimensional work measured flapping:
cells crossing a bare threshold every cycle and the mesh drifting in size.

Three things differ from the two-dimensional case, and the first is not a
detail.

\paragraph{The level table is built by dilation, not by search.}
In two dimensions the distance from every cell to the nearest seed is found by
brute force, $O(n_{\text{cell}} \times n_{\text{seed}})$, and that is affordable.
In three dimensions, 250\,000 cells against thousands of seeds is of order
$10^9$ operations \emph{per remesh}. Instead each seed stamps its neighbourhood
of radius $\lceil \mathrm{thr}_0/h \rceil$, which is
$O(n_{\text{seed}} \times R^3)$.

The dilation is \emph{exact}, not an approximation: the level is monotone
decreasing in distance, so the maximum over seeds is the level of the nearest
seed. That argument is not taken on trust --- the exhaustive version is kept
alongside and the harness requires byte-for-byte equality in five
configurations. Zero cells differ in any of them.

\paragraph{Vorticity stops being a scalar.}
$\omega_z$ suffices in the plane; in space the criterion uses
$|\nabla \times \mathbf{u}|$ with all three components, which costs twelve facet
evaluations per cell against four.

\paragraph{The mesh file format.}
Generic in the dimension, but with six fields per patch instead of four, and
one-based indices. In two dimensions the first version used zero-based indices:
the count of fine cells came out \emph{right} and the positions \emph{wrong},
displaced by one cell of the level, and the measured band fell from 5 to 1.5
cells. A right count with wrong positions appears in no total, so the harness
re-reads the file and checks every patch against the table.

\paragraph{What the loop delivers.}
Measured, base 25 with one level, remeshing every ten steps: the band on the
mesh the solver actually built is 6.03 fine cells against a floor of five; no
kernel support point lands in a cell of different size; no position is lost in
the transfer. The last two must be exactly zero --- the kernel is normalised for
a single $h$, and crossing a level breaks the partition of unity \emph{silently},
because the weights can still sum to nearly one.
